\chapter{KAJIAN PUSTAKA}
\label{chap:kajian-pustaka}
% Target bab: sekitar 10 halaman. Teks dalam lingkungan petunjuk adalah panduan isi; ganti dengan tulisan sendiri.

% ---- Paragraf pembuka (0,25 hlm) ----
\begin{petunjuk}
Jelaskan bahwa bab ini membahas dasar teori dan penelitian terkait, lalu sebut urutan subbab.
\end{petunjuk}

% =====================================================================
\section{\textit{SQL Injection} dan \textit{Taint Analysis}} % 2 hlm
\label{sec:sqli-taint}

\subsection{\textit{SQL Injection} sebagai Kerentanan Aliran Data}
\label{sec:sqli-aliran-data}
\begin{petunjuk}
Mekanisme singkat dan contoh kode Java yang rentan (\texttt{Statement.executeQuery} dengan \textit{string concatenation}), misalnya Listing~\ref{lst:sqli-rentan}.
\end{petunjuk}

\begin{lstlisting}[caption={Contoh kode Java yang rentan terhadap \textit{SQL injection}}, label={lst:sqli-rentan}]
String id = request.getParameter("id");
Statement stmt = conn.createStatement();
ResultSet rs = stmt.executeQuery(
    "SELECT * FROM users WHERE id = '" + id + "'");
\end{lstlisting}

\subsection{Konsep \textit{Taint Analysis}}
\label{sec:konsep-taint}
\begin{petunjuk}
\textit{Source}, \textit{sink}, \textit{sanitizer}, propagasi, dan definisi \textit{taint-flow path}.
\end{petunjuk}

\subsection{Representasi Kode dan Perangkat Analisis}
\label{sec:representasi-kode-tool}
\begin{petunjuk}
AST/CFG/CPG sekilas (hanya yang dipakai \textit{tool}), Joern/CodeQL, dan alasan pemilihan \textit{tool}.
\end{petunjuk}

\subsection{Keterbatasan \textit{Taint Analysis} sebagai \textit{Ground Truth}}
\label{sec:batas-taint}
\begin{petunjuk}
\textit{False positive}/\textit{negative}, dan hanya menangkap sebagian alasan kerentanan.
\end{petunjuk}

% =====================================================================
\section{Deteksi Kerentanan Berbasis \textit{Machine Learning}} % 1 hlm
\label{sec:deteksi-ml}

\subsection{Representasi Kode untuk Model ML}
\label{sec:representasi-ml}
\begin{petunjuk}
\textit{Token} vs graf, dan mana yang dipakai.
\end{petunjuk}

\subsection{Model Detektor}
\label{sec:model-detektor}
\begin{petunjuk}
Model yang dipilih (misal CodeBERT) dan alasannya.
\end{petunjuk}

\subsection{\textit{Dataset} Kerentanan Java}
\label{sec:dataset}
\begin{petunjuk}
Juliet/SARD, OWASP Benchmark, dsb. Label SQLi, keterbatasan, dan metrik performa model (F1, dll.).
\end{petunjuk}

% =====================================================================
\section{\textit{Explainable AI} untuk Model Kode} % 2 hlm
\label{sec:xai}

\subsection{Konsep Dasar XAI}
\label{sec:konsep-xai}
\begin{petunjuk}
\textit{Black box}, \textit{post-hoc} vs \textit{intrinsic}, \textit{local} vs \textit{global}.
\end{petunjuk}

\subsection{SHAP}
\label{sec:shap}
\begin{petunjuk}
Nilai Shapley, cara kerja, kelebihan, dan kelemahan.
\end{petunjuk}

\subsection{LIME}
\label{sec:lime}
\begin{petunjuk}
\textit{Surrogate} lokal berbasis perturbasi, kelebihan, dan kelemahan.
\end{petunjuk}

\subsection{\textit{Attention} sebagai Penjelasan}
\label{sec:attention}
\begin{petunjuk}
Cara kerja di \textit{Transformer} dan perdebatan ``\textit{attention} bukan penjelasan''.
\end{petunjuk}

\subsection{Penjelasan pada Level \textit{Token} untuk Kode}
\label{sec:xai-token}
\begin{petunjuk}
Bagaimana ketiga metode menghasilkan skor per \textit{token} (jembatan ke masalah pemetaan di Subbab~\ref{sec:evaluasi-xai}).
\end{petunjuk}

% =====================================================================
\section{Evaluasi Penjelasan XAI} % 2,5 hlm
\label{sec:evaluasi-xai}

\subsection{\textit{Faithfulness} dan \textit{Plausibility}}
\label{sec:faithfulness}
\begin{petunjuk}
\textit{Comprehensiveness}, \textit{sufficiency}; kenapa penjelasan yang tampak masuk akal belum tentu setia pada model (motivasi riset).
\end{petunjuk}

\subsection{Konsistensi dan Stabilitas}
\label{sec:stabilitas}
\begin{petunjuk}
Kestabilan penjelasan terhadap perubahan input yang tidak mengubah semantik.
\end{petunjuk}

\subsection{Evaluasi Berbasis \textit{Ground Truth}}
\label{sec:eval-ground-truth}
\begin{petunjuk}
ERASER dan \textit{rationale} manusia vs \textit{taint-flow} yang deterministik.
\end{petunjuk}

\subsection{Evaluasi Penjelasan pada Kode}
\label{sec:eval-kode}
\begin{petunjuk}
Level \textit{token}, baris, \textit{statement}; masalah pemetaan skor \textit{token} ke \textit{taint path}.
\end{petunjuk}

\subsection{Keterbatasan Metrik yang Ada untuk Domain Kode}
\label{sec:batas-metrik}

% =====================================================================
\section{Penelitian Terkait dan Celah Penelitian} % 2,25 hlm
\label{sec:penelitian-terkait}

\subsection{XAI untuk Deteksi Kerentanan Kode}
\label{sec:xai-vuln}
\begin{petunjuk}
Karya terdekat: apa yang dijelaskan dan bagaimana dievaluasi.
\end{petunjuk}

\subsection{Tabel Perbandingan Penelitian Terkait}
\label{sec:tabel-perbandingan}
\begin{petunjuk}
Bandingkan 5--8 karya, seperti pada Tabel~\ref{tab:perbandingan}.
\end{petunjuk}

\begin{table}[htbp]
  \caption{Perbandingan penelitian terkait}
  \label{tab:perbandingan}
  \scriptsize
  \setlength{\tabcolsep}{3pt}
  \renewcommand{\arraystretch}{1.2}
  \begin{tabular}{|p{2.2cm}|p{0.9cm}|p{1.9cm}|p{1.3cm}|p{1.5cm}|p{1.9cm}|p{2.5cm}|}
    \hline
    \textbf{Judul} & \textbf{Tahun} & \textbf{Bahasa/\newline Kerentanan} & \textbf{Model} & \textbf{Metode\newline XAI} & \textbf{Metrik\newline Evaluasi} & \textbf{Keterbatasan} \\
    \hline
    & & & & & & \\ \hline
    & & & & & & \\ \hline
    & & & & & & \\ \hline
  \end{tabular}
\end{table}

\subsection{Celah Penelitian}
\label{sec:celah}
\begin{petunjuk}
\begin{enumerate}
  \item Belum ada \textit{ground truth} formal berbasis \textit{taint-flow}.
  \item Belum ada fokus pada Java SQLi.
  \item Belum ada perbandingan SHAP/LIME/\textit{attention} dengan basis yang sama.
\end{enumerate}
\end{petunjuk}

% ---- Ringkasan bab (0,25 hlm) ----
\begin{petunjuk}
Satu paragraf: benang merah antar subbab dan jembatan ke Bab III.
\end{petunjuk}
